% ======================================================================
% Chapter 11 -- Part B: build, deploy, evaluate, business case, end matter
% ======================================================================

\section{Build: How Do We Make It?}
\label{sec:fc-build}

\subsection{The forecast server}

The forecast MCP server is owned by supply-chain analytics and exposes three read-only tools to the exception agent:
\begin{itemize}[itemsep=1pt]
  \item \texttt{get\_forecast(store, product, horizon\_days)} returns the stored quantile forecast from the nightly batch (the 10th to 90th percentiles by day), the model version and the time it was produced. It never runs a model.
  \item \texttt{what\_if(store, product, promotion\_id)} reruns the model for one series with the promotion's discount and display as covariates, and returns the changed quantiles. It is the only tool that triggers inference, and the server counts its calls per night.
  \item \texttt{similar\_items(store, product, k)} returns up to five products with similar sales patterns and their history around comparable events, which is how the agent reasons about new products with little history.
\end{itemize}
The nightly batch itself is not an agent or an MCP tool; it is a scheduled job---a pipeline step in code---that reads the day's sales, runs the model over all series in large batches on accelerators, and writes the forecast store. Separating the batch from the serving tools is what keeps the agent's tools fast and its costs bounded.

\subsection{The exception agent}

The exception agent receives one flagged series at a time with its flag reason and the computed order. Its instruction asks for four things: what is different about this series tonight, what evidence supports an adjustment, a proposed order, and a one-sentence rationale a planner can check in ten seconds. It has ten tool calls per series and no write tools. Its proposal passes a bounds guardrail in code (the adjusted order must be within $\pm$40\% of the computed order and within shelf capacity) and lands in the planner's queue sorted by the money at stake: the difference in expected profit between the computed and the proposed order, estimated from the forecast quantiles.

\subsection{The order calculation}

The newsvendor rule (Equation~\eqref{eq:newsvendor}) is implemented in the replenishment system as a function of each product's price, cost and salvage value, read from the merchandising system nightly. For products with longer shelf life, the rule generalizes to multi-period inventory policies; the replenishment team uses its existing safety-stock logic for those and the newsvendor rule for items with a shelf life of three days or less.

% ----------------------------------------------------------------------
\section{Deploy: How Do We Run It Safely?}
\label{sec:fc-deploy}

\paragraph{Batch on accelerators, serving on Cloud Run.} The nightly batch runs as a scheduled job on accelerator-equipped compute sized to finish in under an hour; the forecast server and the agent run on Cloud Run. Model weights are pinned by version and stored in Artifact Registry, so that the batch uses exactly the validated model.

\paragraph{The nightly timeline as an SLO.} Each stage has a deadline: sales data closed by 10:00~p.m., batch complete by 11:30~p.m., orders computed by midnight, exceptions investigated by 1:00~a.m., planner review by 3:00~a.m., submission by 4:00~a.m. Monitoring alerts on any stage that misses its deadline, and every stage has a fallback: if the batch fails, yesterday's forecasts shifted by one day are used; if the agent fails, exceptions go to planners unannotated; if planners do not act, computed orders are submitted.

\paragraph{Cost controls.} The batch has a fixed size; the variable costs are \texttt{what\_if} calls and agent model calls. Both are budgeted per night, and the agent stops investigating new series when the budget is reached, leaving the rest to planners unannotated.

% ----------------------------------------------------------------------
\section{Evaluate: How Do We Know It Works?}
\label{sec:fc-evaluate}

\subsection{A baseline must be beaten}

\begin{watchout}[Watch out: benchmarks are not your products]
Foundation models for time series are new and improving fast, and published benchmarks do not tell you how a model will do on your products.
\end{watchout} Greenleaf therefore runs a \keyterm{backtest}: for each of the last 52 weeks, it forecasts as if it were that week, using only data available then, and compares the forecasts with what actually sold. Three forecasters are compared: today's moving average, the statistical design's models, and the foundation-model design's foundation model.

Metrics are chosen for the decision the forecast drives. Because orders use quantiles, the main metric is the \emph{pinball loss} at each product's critical ratio, which measures the quality of exactly the quantile that will be ordered. Supporting metrics are a scaled error such as MASE for comparability across products, and, most directly, the \emph{simulated waste and availability} that would have resulted from ordering with each forecaster through the newsvendor rule. The decision rule, written before the backtest is run: the foundation model is adopted if it improves simulated profit on fresh items by at least 2\% over the statistical design's models, with no product category worse by more than 1\%.

\subsection{Exceptions and the agent}

The exception agent is evaluated on last year's flagged series: for each, would its proposed adjustment have improved the outcome compared with the computed order? The metric is the share of adjustments that moved the order in the right direction, weighted by money, and the profit gained or lost in total. Planners' approvals and edits during the pilot provide a second measure: an agent whose proposals are mostly edited is not saving much planner time.

\subsection{Monitoring for drift}

Forecast quality is monitored weekly by category with the same pinball loss, against the backtest's level, and with a simple check that each night's forecast distribution is not wildly different from the previous night's for the same series. A systematic error (\risk{Systematic error})---for example, a data feed that drops a store's sales---shows up first as a sudden shift in many series at once; the batch stage refuses to publish forecasts that move more than a set share of series beyond a threshold, and alerts the analytics team instead.

% ----------------------------------------------------------------------
\section{Business Case: Is It Worth It?}
\label{sec:fc-business}

Greenleaf's fresh sales are about \$310 million a year. Cutting waste from 6.8\% to 4.5\% of fresh sales removes about \$7 million a year of waste at retail value, roughly \$4.6 million at cost; raising availability from 92.5\% to 96\% recovers sales that, at fresh margins, are worth a few million dollars more. The ordering rule alone, applied to better forecasts, produces most of this; the exception agent's main contribution is planner time and better handling of promotions and holidays, where the largest single-day errors occur. Run costs, including the batch, the agent and engineering, are a few hundred thousand dollars a year.

The pilot runs in 10 stores for eight weeks, with 10 matched stores as the comparison, and starts only after the backtest decision. Go if fresh waste falls by at least 1.5 percentage points with availability no lower than in the comparison stores, every nightly deadline is met, and planners accept at least 60\% of the agent's proposals with no or minor edits.

\begin{tcolorbox}[colback=white, colframe=black!60, boxrule=0.6pt, arc=1pt, left=10pt, right=10pt, top=8pt, bottom=8pt, fontupper=\small]
\textbf{To:} Supply-Chain Director \hfill \textbf{From:} Forecasting Design Team\\
\textbf{Re:} Less fresh waste and fuller shelves --- request to backtest and pilot

\smallskip\hrule\smallskip

\textbf{Recommendation.} Approve a 52-week backtest of three forecasting methods, followed by an eight-week pilot in 10 stores using the winner.

\textbf{The problem.} We throw away 6.8\% of fresh sales and are out of stock on fresh items 7.5\% of the time. Planners spend most of their day overriding orders.

\textbf{What changes.} Every product in every store gets a forecast of the range of likely demand each night. Orders are set by each product's economics: for items that spoil, slightly below expected demand. A system investigates the few products where something unusual is happening---promotions, holidays, weather, new lines---and proposes adjustments that planners approve.

\textbf{Value.} Reaching our waste target removes about \$4.6 million a year of waste at cost, with further gains from fewer empty shelves; nightly computing costs are small.

\textbf{Safeguards.} Any adjustment is limited to 40\% of the calculated order. Every stage has a fallback, and orders always reach suppliers on time. The new forecasting model is used only if it beats our best simpler alternative on a year of our own data.
\end{tcolorbox}

% ----------------------------------------------------------------------
\section{Changed Requirement: What If the World Changes?}
\label{sec:fc-changed}

\begin{enumerate}[itemsep=3pt]
  \item \textbf{Intraday reordering.} Greenleaf adds a second, midday delivery for fresh items in its 20 busiest stores. Should forecasting move from batch to on-demand for those stores? Recompute the cost and latency budget.
  \item \textbf{Ten times the series.} An acquisition takes Greenleaf to 400 stores and 12{,}000 products. Which costs scale linearly, which do not, and does the exception share need to change?
  \item \textbf{No accelerators.} Accelerator capacity is unavailable in Greenleaf's region for three months. What do you do on the first night, and for the three months?
\end{enumerate}

% ----------------------------------------------------------------------
\section*{Chapter Summary}
\addcontentsline{toc}{section}{Chapter Summary}

\begin{itemize}[itemsep=2pt]
  \item Batch the predictable and reserve agents for the exceptional: forecasting every series costs about a dollar a night; reviewing every series with an agent would cost millions a year and miss the deadline.
  \item Serve stored results through MCP tools; run inference on demand only for bounded what-if questions.
  \item For perishables, order the quantile at the critical ratio, not the mean; this needs probabilistic forecasts.
  \item A new forecasting model must beat a simple baseline on the company's own backtest, with the decision rule written in advance.
  \item Treat a nightly pipeline's timeline as an SLO with a fallback at every stage.
\end{itemize}

\section*{Review Questions}
\addcontentsline{toc}{section}{Review Questions}

\begin{enumerate}[itemsep=2pt]
  \item Why is the optimal order for the item in Listing~\ref{lst:replenish} below median demand?
  \item Why does the newsvendor rule require a probabilistic forecast?
  \item Why does \texttt{get\_forecast} read stored results while \texttt{what\_if} runs the model?
  \item Explain why an agent reviewing every series is infeasible on both cost and latency.
  \item What is pinball loss, and why is it the right main metric here?
  \item What fallback applies at each stage of the nightly timeline?
\end{enumerate}

\section*{Exercises}
\addcontentsline{toc}{section}{Exercises}

\begin{enumerate}[itemsep=2pt]
  \item \textbf{(Design)} Write the statistical design, with the exception agent added, as a design table, and compare it with the foundation-model design on the trade-off table.
  \item \textbf{(Design)} Write the component card for the exception agent, including its budget behavior.
  \item \textbf{(Build)} In \texttt{replenish.py}, vary the salvage value from \$0 to \$2.40 and plot the critical ratio and the profit gap between ordering the mean and the quantile.
  \item \textbf{(Build)} Add a forecast that is biased 15\% high and show its effect on waste and profit under the newsvendor rule.
\end{enumerate}

\section*{Lab 11: Beat the Baseline}
\addcontentsline{toc}{section}{Lab 11: Beat the Baseline}

\paragraph{Goal.} Backtest a time-series foundation model against statistical baselines on synthetic grocery data, and serve the winner through an MCP server.

\paragraph{Steps.}
\begin{enumerate}[itemsep=2pt]
  \item In \texttt{ch11/}, load two years of synthetic daily sales for 2{,}000 store--product series, with promotions and holidays.
  \item Run a rolling-origin backtest over the last 26 weeks for a moving average, a seasonal exponential-smoothing model and the pinned foundation model. Report pinball loss at each product's critical ratio and simulated profit.
  \item Apply the decision rule written in Section~\ref{sec:fc-evaluate}.
  \item Write the winning forecasts to a store and serve them through \texttt{get\_forecast}; implement \texttt{what\_if} for the foundation model.
  \item Measure batch throughput on your hardware and replace the assumed numbers in \texttt{replenish.py}.
\end{enumerate}

\paragraph{Deliverable.} The backtest table, the decision, and the corrected compute budget.

\section*{Further Reading}
\addcontentsline{toc}{section}{Further Reading}

\begin{itemize}[itemsep=2pt]
  \item \citet{silver2017inventory} covers the newsvendor problem and the multi-period policies used for longer-life products.
  \item \citet{ansari2024chronos} and \citet{das2024timesfm} describe two influential time-series foundation models; read their evaluation sections with this chapter's backtest discipline in mind.
\end{itemize}
